\documentclass[11pt,a4paper]{article}
\usepackage[margin=1in]{geometry}
\usepackage{hyperref}
\usepackage{enumitem}
\usepackage{graphicx}
\usepackage{xcolor}

% -----------------------------------------------------
% bvraghav-tiet-ucs503p-article.sty
% -----------------------------------------------------

% Variable: Lab Instructor
\makeatletter \newcommand{\BvrSetLabInstructor}[1]{%
  \gdef\Bvr@LabInstructor{#1}}
% default
\newcommand{\Bvr@LabInstructor}{Dr. Raghav %
  B. Venkataramaiyer}
% public accessor
\newcommand{\BvrLabInstructorValue}{\Bvr@LabInstructor}
\makeatother

% Add subtitle
\usepackage{titling}
% -- Set title bold
\pretitle{\begin{center}\bfseries\fontsize{18bp}{18bp}\selectfont}
\makeatletter
\newcommand{\BvrSubtitle}[1]{%
  \posttitle{%
    \par\end{center}
    \begin{center}\large#1\end{center}
    \vskip 0.5em}%
}
\makeatother

% Hints in the section title(s)
\newcommand{\BvrHint}[1]{%
  {\normalfont\normalsize\itshape\hfill #1}}

% -----------------------------------------------------

\title{CampusConnect: A Unified Campus Event and
  Society Management Platform}

\BvrSetLabInstructor{Dr. Raghav B. Venkataramaiyer}

\BvrSubtitle{%
  Project Proposal for UCS503P  \\
  Submitted to: \BvrLabInstructorValue \\ \bfseries
  Thapar Institute of Engineering and Technology%
}

\author{
  Nitin Kumar, CSED%
  \thanks{Roll No: \texttt{1024030339}, %
    Email: \texttt{nkumar\_be24\@thapar.edu}}
  \and 
  Mohan Arora, CSED%
  \thanks{Roll No: \texttt{1024031087}, %
    Email: \texttt{marora\_be24\@thapar.edu}}
  \and 
  Prateek Malpani, CSED%
  \thanks{Roll No: \texttt{1024030340}, %
    Email: \texttt{pmalpani\_be24\@thapar.edu}}
}

\date{\today}

\begin{document}
\maketitle

\section{Higher-order goal%
  \BvrHint{\ldots secondary framing}}
This project aims to improve campus governance by
reducing the time and effort required to take a society
event from an idea to a registerable campus-wide event.
While the topic is broader than any single event,
the immediate engineering objective is to translate
\emph{efficient event administration} into a measured
software solution that connects societies, the
administration, and students on one platform.

\section{Time-to-value %
\BvrHint{\ldots primary heuristic}}
\begin{itemize}[leftmargin=0.7cm]
    \item We will deliver meaningful increments quickly by prototyping the core user workflow and validating it with actual campus roles within a short iteration window.
    \item We will prioritize fast feedback over perfection by using weekly iterations (and targeting CI-backed automation early) to reduce release friction.
    \item Success is defined by what can be delivered and measured in the first iteration, then improved based on user feedback.
\end{itemize}

\section{Problem Statement}
Managing the activities of many student societies
across a large campus is fragmented and slow.
Common pain points include:
\begin{itemize}[leftmargin=0.7cm]
    \item \textbf{Fragmentation:} proposals, approvals
      and announcements live in emails, notice boards
      and paper forms, causing lost context.
    \item \textbf{Low transparency:} students cannot
      easily see what is happening across campus or
      whether an event has been approved and will run.
    \item \textbf{Slow coordination:} an authority has
      no structured queue to review, approve, reject or
      send back an event proposal with clear reasons.
    \item \textbf{Poor data quality:} registrations are
      manual, capacity is untracked, and there is no
      reliable record of who organizes or attends an
      event.
\end{itemize}

\textbf{Why this matters now (contextual relevance):} On
a campus with dozens of active societies, even small
process improvements can noticeably reduce the delay
between an event idea and a full room of participating
students.

\section{Proposed Solution %
\BvrHint{\ldots Software Proposition}}
\subsection{Overview}
Build a \textbf{web-based} ``CampusConnect'' platform with
the following components:
\begin{itemize}[leftmargin=0.7cm]
    \item \textbf{Society portal:} draft and submit event
      proposals with date, venue, budget, objectives and
      coordinators.
    \item \textbf{Admin review queue:} approve, reject or
      request changes on proposals with comments;
      manage categories, societies, users and activity.
    \item \textbf{Public campus index:} a login-free
      catalogue of upcoming events, categories, and a
      complete society directory with profiles.
    \item \textbf{Student module:} register and cancel
      for events, save events, and view registration
      history.
    \item \textbf{Notifications:} in-app notifications at
      every lifecycle step (submitted, approved,
      rejected, changes requested, event cancelled).
\end{itemize}

\subsection{Core workflow}
\begin{enumerate}[leftmargin=0.7cm]
    \item A society drafts a proposal and submits it;
      the proposal enters the admin review queue.
    \item Admin approves (an event is created and
      published), requests changes (the society revises
      and resubmits), or rejects with a reason.
    \item The approved event appears on the public index,
      under its category, and on the society's public
      profile.
    \item Students browse freely and register when ready;
      capacity and live registration counts are tracked.
\end{enumerate}

\subsection{Operational constraints for an educational, campus setting}
\begin{itemize}[leftmargin=0.7cm]
    \item Keep the solution usable on low-to-mid range
      devices via responsive web design.
    \item Avoid dependence on advanced research
      algorithms; focus on workflow, correctness, and
      system hardening.
    \item Seed the system with the real TIET society list
      (57 societies across 11 categories) so it is
      immediately meaningful to evaluators.
\end{itemize}

\section{Solution Approach %
\BvrHint{\ldots Engineering focus}}
This is an engineering course project, so we focus on
workflow automation and system correctness rather than
novel algorithm research.

\subsection{Web architecture %
\BvrHint{\ldots web-based solution}}
A typical 3-tier web architecture:
\begin{itemize}[leftmargin=0.7cm]
    \item \textbf{Frontend:} React 19 + TypeScript + Vite
      SPA with role-based layouts (public, student,
      society, admin).
    \item \textbf{Backend API:} Node.js + Express +
      TypeScript; modules for auth, public, student,
      society and admin routes; zod validation on every
      route.
    \item \textbf{Database:} SQLite via Prisma ORM
      (PostgreSQL-ready), storing users, societies,
      members, proposals, events, registrations,
      notifications and activity logs.
\end{itemize}

\subsection{Security and correctness %
\BvrHint{\ldots non-research}}
\begin{itemize}[leftmargin=0.7cm]
    \item JWT access + refresh tokens with role-based
      access control (\texttt{ADMIN}, \texttt{SOCIETY},
      \texttt{USER}).
    \item Bcrypt password hashing; helmet, rate limiting
      and restricted CORS; uploads restricted by MIME
      whitelist and size cap.
    \item Status-transition guards so a proposal cannot
      be approved twice or an event cancelled twice.
\end{itemize}

\subsection{CI/CD and fast delivery enabler}
CI/CD will be used to:
\begin{itemize}[leftmargin=0.7cm]
    \item Automatically build and run tests on every
      change.
    \item Produce repeatable documentation builds
      (mkdocs deployed to GitHub Pages on every push to
      \texttt{master}/\texttt{main}).
    \item Enable safe and frequent releases of
      incremental features.
\end{itemize}

\section{Evaluation Criterion %
\BvrHint{\ldots measurable and attributable}}
We define success using measurable metrics tied to the
core workflow.

\subsection{Primary evaluation metric}
\textbf{Proposal-to-publication latency:} time between a
society submitting a proposal and the administration
publishing the resulting event.
\begin{itemize}[leftmargin=0.7cm]
    \item Target: median latency $\le$ 48 hours during the
      pilot window.
    \item Attribution: measured from database timestamps
      (\texttt{submittedAt} vs.\ \texttt{reviewedAt} and
      \texttt{publishedAt}).
\end{itemize}

\subsection{Secondary evaluation metrics}
\begin{itemize}[leftmargin=0.7cm]
    \item \textbf{Review completeness:} percentage of
      submitted proposals resolved (approved or
      rejected) without rework requests.
    \item \textbf{Registration accuracy:} capacity never
      oversubscribed beyond the configured cap, and
      attendance matches registrations.
    \item \textbf{Discovery coverage:} every approved
      event is reachable from the public index within one
      click of its category page.
    \item \textbf{Reliability:} the integration test
      suite passes end-to-end (\texttt{npm test} in the
      backend and \texttt{tsc -b} + \texttt{vite build}
      in the frontend).
\end{itemize}

\subsection{Pilot validation plan %
\BvrHint{\ldots high-level}}
\begin{itemize}[leftmargin=0.7cm]
    \item Recruit 10--20 pilot users (students) and 3--5
      roles from society and administration (simulated if
      real staff are unavailable).
    \item Compare metrics within the iteration window by
      logging baseline estimates via short interviews or
      existing process observations.
\end{itemize}

\section{Scalability %
\BvrHint{\ldots with foresight}}
The proposition should scale in both theoretical and
practical senses.

\begin{enumerate}[leftmargin=0.7cm]
    \item \textbf{Scalable (theoretically):} The system
      should support growth in the number of societies,
      events and concurrent student traffic by:
    \begin{itemize}[leftmargin=0.7cm]
        \item decoupling components (frontend/backend),
        \item enabling pagination and indexing for common
          queries (events by date/status, societies by
          category),
        \item using stateless API design.
    \end{itemize}

    \item \textbf{Operationally deployable (practically):}
      The system should run on common web hosting:
    \begin{itemize}[leftmargin=0.7cm]
        \item managed database (PostgreSQL via Prisma),
        \item containerized or platform-hosted
          deployment,
        \item CI/CD pipeline for staging/production.
    \end{itemize}
\end{enumerate}

\section{Engine availability heuristic}
A mature set of ``engines'' (tooling/services) is
available to implement this as a web-based project:
\begin{itemize}[leftmargin=0.7cm]
    \item React + Vite for the SPA, and Express for REST
      APIs.
    \item Prisma ORM with a managed relational database.
    \item Token-based authentication (JWT).
    \item zod for validation; multer for uploads.
    \item Vitest + supertest for unit/integration tests.
    \item CI runner and GitHub Pages deployment target.
\end{itemize}
We will use these mature components to focus on workflow
design, correctness, and measurement rather than
inventing new infrastructure.

\section{Project scope and deliverables %
\BvrHint{\ldots iteration-friendly}}
\subsection{Initial deliverable %
\BvrHint{\ldots first iteration}}
\begin{itemize}[leftmargin=0.7cm]
    \item Society proposal form with full validation and
      submission
    \item Admin review queue: approve / reject / request
      changes
    \item Public event index and event detail pages
    \item Student registration with capacity tracking
    \item CI: build + backend integration tests + linting
\end{itemize}

\subsection{Subsequent deliverables}
\begin{itemize}[leftmargin=0.7cm]
    \item Society directory with per-society profiles and
      event lists
    \item Notification layer at every lifecycle step
    \item Admin dashboard for statistics and activity log
    \item Saved events, registration history, and
      search/filtering improvements
\end{itemize}

\section{Risks and mitigations}
\begin{itemize}[leftmargin=0.7cm]
    \item \textbf{Data privacy / consent:} minimize
      collected data; expose no password or internal
      identifiers through public APIs (test-verified).
    \item \textbf{Adoption:} provide an editorial, public
      discovery experience so the read-side needs zero
      training; keep role dashboards minimal and guided.
    \item \textbf{Evaluation measurement ambiguity:}
      instrument timestamps (\texttt{submittedAt},
      \texttt{reviewedAt}, \texttt{publishedAt}) and define
      clear status transitions before development begins.
    \item \textbf{Operational issues in pilot:} use CI
      early; ensure repeatable \texttt{npm run} scripts
      for seed, test, and build.
\end{itemize}

\section{Summary}
This project proposes a deployable web platform that
takes society events from proposal to registration on one
campus index.  It meets the course criteria by
emphasizing time-to-value through rapid weekly
iterations, implementing CI/CD to support frequent safe
releases, and using measurable evaluation criteria
(especially proposal-to-publication latency).  It remains
focused on engineering workflow and scalability readiness
rather than research-driven novelty.

\end{document}